\documentclass[12pt]{article}
\usepackage[utf8]{inputenc}
\usepackage[T1]{fontenc}
\usepackage{hyperref}
\title{U Language Resources for K-12 Curriculum Development}
\author{Compiled by Claude Sonnet 5 (claude-sonnet-5) for Andr\'as Kornai}
\date{}
\begin{document}
\maketitle

\section{Introduction}
U (ISO 639-3 uuu) is an Austroasiatic language of the Palaungic branch, spoken by roughly 40,000 people in Yunnan Province, China, and possibly in Myanmar. Its speakers are officially classed as Bulang in China. It should not be confused with M\h{a}ng (zng) or Kh\'ang (kjm), which are separate languages. The search found only encyclopedic and catalogue entries. It found no curriculum material, primers, or school programmes.

\section{Primary Educational Resources}
No primers, textbooks, or school materials for U turned up in the searches.

\section{Linguistic and Reference Materials}
\begin{itemize}
\item Wikipedia overview of the U language: \url{https://en.wikipedia.org/wiki/U_language}
\item Glottolog language page (the result was listed under the identifier mang1378 and may be a different Mang-named entry, so check it before relying on it): \url{https://glottolog.org/resource/languoid/id/mang1378}
\item Endangered Languages Project bibliography page (entry 1236; the search result did not show which language this entry is, so verify): \url{https://www.endangeredlanguages.com/lang/1236/bibliography?hl=fr}
\end{itemize}

\section{Religious and Community Resources}
None found.

\section{Digital and Audio Resources}
None found for U itself. The Endangered Languages Project lists a samples page for entry 1236 whose language was not verified: \url{https://endangeredlanguages.com/lang/1236/samples/10372?hl=fr}

\section{Curriculum Development Context}
Very little is documented about U. It is not known to be taught in schools, and direct evidence of its use is lacking. Any curriculum work would have to begin with fieldwork and community consultation, and with confirming exactly which speech variety is meant.
\end{document}